\documentclass[11pt,a4paper]{article}

\usepackage[utf8]{inputenc}
\usepackage[T1]{fontenc}
\usepackage[ngerman]{babel}
\usepackage[a4paper,margin=2.5cm]{geometry}
\usepackage{microtype}
\usepackage{graphicx}
\usepackage{booktabs}
\usepackage{tabularx}
\usepackage{enumitem}
\usepackage[hidelinks]{hyperref}
\usepackage{xurl}

\setlist{nosep,leftmargin=*}
\setlength{\parindent}{0pt}
\setlength{\parskip}{0.55em}
\setlength{\emergencystretch}{2em}
\graphicspath{{../../docs/daten/img/}{../../docs/modelle/img/}{../figures/}}

\title{Abschlussbericht\\Kulturpflanzen- und Entwicklungsstadiumsklassifikation}
\author{Dom\"anenprojekt 2\\Team: AE-ROOT}
\date{Stand: \today}

\begin{document}
\maketitle
% Teammitglieder, Abgabedatum und den ausgewerteten Code-Stand erg\"anzen.
% Ergebnisse mit MLflow-Run-IDs und reproduzierbaren Artefakten belegen.
% Haupttext knapp halten; Detailtabellen und zus\"atzliche Plots in den Anhang.

\tableofcontents
\clearpage

\section{Ziel und Dom\"ane}
% Die Aufgabe erl\"autern: Crop und Stage aus EO-1-Hyperion-Spektren vorhersagen.
% Fachliche Motivation und vorgesehenen Einsatz beschreiben und abgrenzen.
% Zentrale Herausforderungen nennen: ungleiche Klassen, \"ahnliche Spektren,
% wenige Beispiele f\"ur manche Crop-Stage-Kombinationen.
% Untersuchungsfragen formulieren: Modellstrategie, Nutzen der Metadaten und
% Trade-off zwischen Datenreduktion und Klassifikationsleistung.

\section{Daten und explorative Analyse}

\subsection{Datengrundlage und Zielgr\"o\ss en}
% Herkunft, Umfang, Spektralbereich und Bandanzahl anhand der Daten belegen.
% AEZ, Month, Crop und Stage erkl\"aren; bekannte Herkunftsl\"ucken offenlegen.
% Gelabelte Entwicklungsdaten und ungelabelte Abgabedaten klar unterscheiden.
% Beobachtete Kombinationen nicht mit fachlich m\"oglichen gleichsetzen.

\subsection{Zentrale Befunde und Konsequenzen}
% Klassenverteilung und seltene Crop-Stage-Paare mit Fallzahlen zeigen.
% Fehlende Werte, leere B\"ander, Duplikate und auff\"allige Spektren beschreiben.
% Spektrale \"Uberlappungen, Bandkorrelationen sowie Zusammenh\"ange mit AEZ und
% Month auf die Befunde beschr\"anken, die Modellentscheidungen begr\"unden.
% Wenige aussagekr\"aftige Abbildungen ausw\"ahlen und im Text interpretieren.

\section{Methodik und Validierung}

\subsection{Datenaufbereitung und Merkmalsbildung}
% Gew\"ahlte Bereinigung, Imputation und Skalierung fachlich begr\"unden.
% Verwendete Bandselektion, Vegetationsindizes, PCA oder weitere Merkmale
% beschreiben; genutzte Metadaten und deren Kodierung nennen.
% Dokumentieren, welche Transformationen ausschlie\ss lich auf den jeweiligen
% Trainingsanteilen innerhalb der gemeinsamen Pipeline angepasst werden.

\subsection{Evaluationsprotokoll und Reproduzierbarkeit}
% Gemeinsame Split- und Fold-Artefakte, Anteile, Stratifizierung und Seed nennen.
% Cross-Validation f\"ur die Hyperparametersuche von der anschlie\ss enden
% Holdout-Auswertung trennen; tats\"achliches Auswahlverfahren offenlegen.
% Falls wiederholt Holdout-Ergebnisse verglichen wurden, die dadurch begrenzte
% Unabh\"angigkeit der finalen Bewertung diskutieren.
% Balanced Accuracy als Hauptmetrik sowie Macro-F1 f\"ur Crop und Stage angeben.
% Samples-F1 f\"ur die gemeinsame Crop-/Stage-Vorhersage pr\"azise definieren.
% Umgang mit Klassenungleichgewicht und ung\"ultigen Vorhersagepaaren erl\"autern.
% Code-Commit, Umgebung (uv.lock), Konfigurationen, MLflow-Run-IDs und Schritte
% zur Reproduktion bzw. Vorhersage auf den Abgabedaten referenzieren.

\section{Modellvergleich und Auswahl}

\subsection{Baseline und untersuchte Strategien}
% Baseline und tats\"achlich untersuchte Strategien beschreiben, beispielsweise
% gemeinsame Klassen, hierarchische Klassifikation und Multi-Task-Learning.
% Relevante Modellfamilien mit ihren Vor- und Nachteilen begr\"unden.
% Suchr\"aume, Suchbudget und wichtigste gew\"ahlte Hyperparameter zusammenfassen.
% Nur entscheidungsrelevante Experimente aufnehmen; vollst\"andige Laufhistorie
% \"uber Experimentdokumentation und MLflow referenzieren.

\subsection{Vergleich und finales Modell}
% Modelle unter vergleichbaren Splits, Metriken und Eingaben gegen\"uberstellen.
% Unterschiede in Suchbudget oder Evaluationsbedingungen kenntlich machen.
% Eine kompakte Ergebnistabelle mit Baseline und wichtigsten Kandidaten zeigen.
% CV-Ergebnisse und Holdout-Ergebnisse getrennt kennzeichnen; vorhandene
% Streuungen berichten, keine Unsicherheiten ohne entsprechende Messung angeben.
% Modellwahl anhand der vorab benannten Kriterien begr\"unden; Leistung,
% Komplexit\"at, Rechenaufwand und Robustheit gegeneinander abw\"agen.

\section{Datenreduktion}
% Vollst\"andigen Input und ausgew\"ahlte reduzierte Varianten vergleichen.
% Anzahl der Eingabeb\"ander, abgeleiteten Features und PCA-Komponenten klar
% unterscheiden: weniger Komponenten bedeuten nicht weniger ben\"otigte B\"ander.
% Auswahlverfahren und Einbettung in die Trainingspipeline dokumentieren.
% Leistung gegen Eingabeumfang als Tabelle oder Plot zeigen, mit gleicher
% Evaluationsgrundlage; gegebenenfalls Laufzeit und Speicherbedarf erg\"anzen.
% Den Trade-off bewerten: Welche Reduktion ist bei welchem Leistungsverlust
% vertretbar? Schlussfolgerungen auf die tats\"achlich gepr\"uften Varianten begrenzen.

\section{Ergebnisse und Interpretation}

\subsection{Vorhersageleistung und Fehleranalyse}
% Finale Kennzahlen f\"ur Crop und Stage sowie gemeinsame Samples-F1 berichten.
% Verbesserungen zur Baseline quantifizieren; Datenbasis jeder Kennzahl nennen.
% Confusion Matrices sowie Precision, Recall und Fallzahlen je Klasse zeigen.
% H\"aufige Verwechslungen und Schw\"achen bei seltenen Klassen diskutieren.
% Bei hierarchischen Modellen m\"ogliche Fehlerfortpflanzung betrachten.
% F\"ur ungelabelte Abgabedaten keine selbst berechneten G\"utema\ss e behaupten;
% externe Bewertung nur bei tats\"achlich vorliegender R\"uckmeldung aufnehmen.

\subsection{Fachliche Einordnung und Erkl\"arbarkeit}
% Ergebnisse mit spektralen Mustern und Dom\"anenwissen in Beziehung setzen.
% Gemessene Befunde von plausiblen, aber ungepr\"uften Erkl\"arungen trennen.
% Falls untersucht: Feature Importance, relevante Wellenl\"angen und Nutzen
% von AEZ/Month erl\"autern; Verfahren und Grenzen der Erkl\"arung nennen.
% Korrelation nicht als Kausalit\"at auslegen; bei korrelierten B\"andern die
% eingeschr\"ankte Interpretierbarkeit einzelner Wichtigkeiten diskutieren.

\section{Modellkarte}
% Als kompakten Steckbrief des tats\"achlich finalen Modells ausarbeiten.
% Falls eine separate Modellkarte verlangt wird, diesen Inhalt auslagern.

\subsection{Identit\"at und vorgesehene Nutzung}
% Modellname, Version, Code-Commit, MLflow-Run-ID und Modellartefakt nennen.
% Zielgruppe, Verwendungszweck und nicht abgedeckte Anwendungen festhalten.

\subsection{Ein- und Ausgaben, Daten und Leistung}
% Ben\"otigte Spektralb\"ander, Einheiten, Metadaten und Vorverarbeitung angeben.
% Ausgabelabels und den Umgang mit Crop-Stage-Kombinationen beschreiben.
% Trainingsbasis, Evaluationsverfahren und zentrale Kennzahlen zusammenfassen
% oder eindeutig auf die entsprechenden Tabellen im Bericht verweisen.

\subsection{Annahmen, Grenzen und Risiken}
% Annahmen zu Sensor, Bandbelegung, Messbedingungen und Datenverteilung nennen.
% Grenzen f\"ur seltene Klassen, unbekannte Kombinationen, fehlende Eingaben und
% abweichende Regionen, Aufnahmezeitpunkte oder Sensoren dokumentieren.
% M\"ogliche r\"aumliche/zeitliche Abh\"angigkeiten und deren Pr\"ufbarkeit benennen;
% ein zuf\"alliger Split belegt keine \"Ubertragbarkeit auf neue Regionen.
% Abh\"angigkeit von AEZ/Month und m\"ogliche Verteilungsverschiebungen diskutieren.
% Nicht vorhandene Absicherungen wie Kalibrierung, Ablehnung unbekannter Daten
% oder Praxistests ausdr\"ucklich als nicht gepr\"uft kennzeichnen.

\section{Fazit und Ausblick}
% Untersuchungsfragen beantworten: bestes begr\"undet gew\"ahltes Modell,
% Zielerreichung und Ergebnis der Datenreduktionsstudie knapp zusammenfassen.
% St\"arkste Erkenntnisse und wichtigste verbleibende Grenzen nennen.
% Wenige konkrete n\"achste Schritte aus den beobachteten Schw\"achen ableiten.

\begin{thebibliography}{9}
% Verwendete Fachliteratur, Datensatz- und Methodenquellen mit \bibitem aufnehmen
% und an den passenden Stellen mit \cite referenzieren.
% Beispielstruktur: \bibitem{schluessel} Autor. Titel. Jahr. DOI oder URL.
\end{thebibliography}

\appendix
\section{Erg\"anzende Ergebnisse und Reproduktion}
% Nur bei Bedarf: vollst\"andige Klassentabellen, weitere Confusion Matrices,
% Suchr\"aume, Bandlisten und ausf\"uhrliche Reproduktionsbefehle aufnehmen.
% Jede zus\"atzliche Tabelle oder Abbildung aus dem Haupttext referenzieren.
% Separate Abgabebestandteile nicht vergessen: Vorhersage-CSV,
% Arbeitszeitnachweise und je Teammitglied 0,5--1 Seite eigener Beitrag.

\end{document}